\documentclass[11pt,a4paper]{article}
\usepackage[utf8]{inputenc}
\usepackage{graphicx}
\usepackage{amsmath}
\usepackage{geometry}
\usepackage{float}
\geometry{a4paper, margin=2.5cm}

\title{DD2437 Lab Assignment 1a: \\ Classification with a Single-Layer Perceptron}
\author{Cardaccia Livia, Felbor Tytus, Prete Giovanni}
\date{\today}

\begin{document}

\maketitle

\section{Main objectives and scope of the assignment}
The primary objective of this assignment is to investigate the capabilities and limitations of single-layer perceptrons (Threshold Logic Units) for binary classification tasks. Specifically, the study compares the classical Perceptron learning rule with the Delta (Widrow-Hoff) rule, evaluating their convergence dynamics, sensitivity to hyperparameters, and behavior in both sequential (online) and batch learning modes. 

Furthermore, the scope includes analyzing the geometric limitations of linear classifiers. This is achieved by evaluating the network's performance on datasets that are not linearly separable (overlapping and bimodal distributions), investigating the impact of the bias weight on the decision boundary, and observing how sampling bias (via data subsampling) skews the network's generalization.

\section{Methods}
\subsection{Data Generation and Representation}
Datasets were generated using 2D multivariate normal distributions. For linearly separable data, 100 samples per class were generated with means $\mu_A = [1.0, 0.5]$ and $\mu_B = [-1.0, 0.0]$, using a standard deviation of $\sigma = 0.3$. For overlapping (non-linearly separable) data, the standard deviation was increased to $\sigma = 0.5$. A third dataset featured a bimodal distribution for Class A. To enable the network to learn an offset decision boundary, the input matrix $X$ was augmented with a row of ones (the bias trick). Target labels $T$ were set to $\{0, 1\}$ for the Perceptron rule and converted to $\{-1, 1\}$ for the Delta rule to ensure correct gradient direction.

\subsection{Learning Algorithms}
Three algorithmic variants were implemented from scratch using NumPy:
\begin{itemize}
    \item \textbf{Perceptron Rule (Sequential):} Updates weights sample-by-sample based on the error computed \textit{after} the step activation function ($\Delta W = \eta (t - y) x^T$).
    \item \textbf{Delta Rule (Sequential):} Updates weights sample-by-sample using the continuous error measured \textit{before} thresholding ($\Delta W = \eta (t - Wx) x^T$).
    \item \textbf{Delta Rule (Batch):} Computes the error across the entire dataset simultaneously. To allow direct comparison using the same learning rate ($\eta = 0.01$), the accumulated gradient was normalized by the number of samples $N$ ($\Delta W = \frac{\eta}{N} (T - WX) X^T$).
\end{itemize}
Performance was tracked per epoch using Mean Squared Error (MSE) and the misclassification rate. Weights were initialized from a zero-mean normal distribution.

\section{Results and discussion}

\subsection{Classification with a single-layer perceptron}

\subsubsection{Algorithm Comparison and Convergence}
The algorithms were first evaluated on the linearly separable dataset. As shown in Figure \ref{fig:separable_comparison}, both the Delta Sequential and Delta Batch algorithms successfully converged, reaching a misclassification rate of 0 within 15 epochs. The MSE for the Batch mode decayed smoothly and monotonically, confirming a stable descent along the error surface. In contrast, the Sequential mode exhibited slight local oscillations before converging, which is characteristic of the stochastic nature of per-sample updates. Random weight initialization affected the exact starting error and slightly shifted the number of epochs to convergence, but did not prevent either variant from finding a valid separating hyperplane.

\begin{figure}[H]
    \centering
    \includegraphics[width=0.8\textwidth]{metrics_comparison_test1.png}
    \caption{Learning curves comparing the Delta Rule (Sequential) and Delta Rule (Batch) on linearly separable data. Both variants successfully reduce the misclassification rate to zero.}
    \label{fig:separable_comparison}
\end{figure}

\subsubsection{Impact of Removing the Bias}
When the bias term was removed ($w_0 = 0$), the decision boundary was geometrically constrained to pass exactly through the origin ($w_1x_1 + w_2x_2 = 0$). Experimental verification confirmed the hypothesis: the bias-free perceptron only converged when the distribution means $\mu_A$ and $\mu_B$ were located in strictly opposite quadrants (e.g., upper-right and lower-left). When both means were placed in the same quadrant, the network failed to separate the data entirely, resulting in a strictly positive misclassification rate [INSERT FINAL MISCLASSIFICATIONS HERE], as a line anchored at the origin could not physically pass between the clusters.

\subsection{Classification of data that are not linearly separable}

\subsubsection{Overlapping Distributions}
For the dataset with increased variance ($\sigma=0.5$), the classes overlapped, making perfect linear separation impossible. Figure \ref{fig:overlapping} demonstrates the failure mode of the single-layer perceptron. Both the classical Perceptron and the Delta Rule failed to reach zero misclassifications, heavily oscillating or plateauing at a minimum error rate of approximately 2\%. Notably, the Delta rule achieved a lower and more stable MSE, as it optimizes for the continuous distance to the target rather than solely focusing on hard misclassifications.

\begin{figure}[H]
    \centering
    \includegraphics[width=0.8\textwidth]{metrics_comparison_test3.png}
    \caption{Learning curves for Perceptron vs. Delta Rule on overlapping, non-linearly separable data. Neither algorithm reaches zero misclassifications.}
    \label{fig:overlapping}
\end{figure}

\subsubsection{Bimodal Data and Sampling Bias}
The network was subsequently tested on a non-linearly separable bimodal dataset under four subsampling conditions. Subsampling heavily skewed the decision boundary. For instance, removing 50\% of Class A caused the network to prioritize minimizing errors on the majority class (Class B), shifting the decision boundary toward Class A and degrading the True Positive Rate [INSERT PER CLASS ACCURACY HERE].

As observed in the subsampling plots (e.g., \texttt{removed\_25\_A\_25\_B.png}), the sequential algorithms (Perceptron and Delta) exhibited severe thrashing, oscillating around a 20--30\% misclassification rate. Interestingly, the Delta Batch algorithm stabilized but converged to a highly suboptimal local minimum, resulting in a persistent misclassification rate of roughly 65\%. This suggests that for highly asymmetric, uncentered data, the averaged batch gradient directed the network toward a local minimum where it likely collapsed to predicting the majority class, whereas the stochastic noise of the sequential updates allowed the online rules to escape and achieve better, albeit unstable, empirical accuracy.

\section{Final remarks}
This lab demonstrated the fundamental mechanics and theoretical limitations of the single-layer perceptron. While both the Perceptron and Delta rules successfully classify linearly separable data, the Delta rule's optimization of continuous error yields smoother gradients and greater stability. The implementation of batch learning highlighted the necessity of gradient normalization to prevent exploding weights and revealed its susceptibility to local minima in asymmetric, non-separable datasets. Finally, the subsampling experiments emphasized that a model's generalization is only as robust as the representation of its training data; sampling bias physically rotates the decision boundary, artificially inflating accuracy for majority classes at the expense of minority classes.

\end{document}